\section{Introduction}

Sign-language translation systems must solve two different problems: extract
linguistic evidence from motion, and turn that evidence into fluent text.  Large
language models (LLMs) are strong candidates for the second problem, but a visual
encoder still needs a well-defined interface through which to communicate its
prediction.  Continuous visual prompts and embedding alignment offer one route;
discrete language-model tokens offer another.

The first published version of \system{} mapped 2D body keypoints to the token
embedding space of a frozen LLM.  Lightweight temporal convolutions and a
Transformer encoded the video; learned cross-attention queries then imitated a
fixed-length sequence of LLaMA embeddings.  It reported a preliminary combined
MSE--cosine validation loss near $8\times10^{-4}$
\citep{velasquez2026imitator}.  This was attractive because the LLM could consume
visual vectors without a conventional gloss-recognition stage.  However, a small
geometric loss does not imply that an approximate vector decodes to the intended
token.  Contextual embedding spaces are anisotropic, and proximity alone is not a
transcription metric \citep{ethayarajh2019contextual}.  Moreover, the original
partitions were not signer-exclusive, so that study did not establish
generalization to an unseen person.

We retain the central idea---a compact \emph{Imitator} learns the input language of
a frozen LLM---but change what is imitated.  The current model predicts the exact
token IDs produced by the \Gemma{} tokenizer.  The output can be decoded directly,
measured with sequence metrics, or passed through the ordinary text interface of
\Gemma{} 3n \citep{gemma3n}.  The intended system is consequently modular:
\system{} is responsible for visually grounded token transcription, while \Gemma{}
is responsible for linguistic realization and correction only to the extent that
the predicted text contains enough context.

This distinction matters empirically.  Our current data are synthetic streams
constructed from isolated signs in random order.  They are useful for testing
whether the visual model can produce ordered token IDs without exploiting a
language prior, but they are almost the worst possible setting for LLM correction:
if the only word is replaced, no linguistic evidence remains.  We therefore report
both the upstream transcription result and a downstream corruption audit, including
the negative outcome.

Our contributions are:
\begin{enumerate}
    \item We reformulate \system{} from continuous embedding imitation as direct,
          autoregressive prediction of exact \Gemma{} token IDs.
    \item We evaluate this interface under a signer-disjoint LOSO protocol, with a
          configuration frozen on development folds and prospectively tested on two
          outer signers.
    \item We measure the downstream \Gemma{} interface rather than assuming that an
          LLM will repair upstream mistakes.  The measured error budget is zero at
          the tested nonzero corruption levels for isolated labels, motivating a
          future sentence-level test with genuine linguistic redundancy.
    \item We retain temporal-order restoration and segment permutation as a causal
          architectural validation, but treat them as enabling evidence for
          \system{} rather than as the paper's central task.
\end{enumerate}
